\subsection{Large MuDst sample: method 2 (overlay) and comparison}
\label{sec:bigm2}

Method 2 (\fpath{StTofMEOverlayMaker} v3, mode 1: partner hits and partner global tracks overlaid on this event, production
\fpath{StBTofMatchMaker}/\fpath{StBTofCalibMaker} re-run; Sec.~\ref{sec:mudst}) was run on the same 102 files, one job per file
(\fpath{scripts/ovbig.sub}). The outputs are merged into \fpath{data/ov_big.root} (\fpath{scripts/merge_ovbig.sh}):
\ovbigEvAll{} events, \ovbigEvMtwo{} used for $m^2$, on average \ovbigNTrkAdd{} partner global tracks added per overlay. The analysis
chain is the one of Sec.~\ref{sec:big}: \fpath{scripts/make_ov.C} (tag \fcmd{big}), \fpath{scripts/make_qm2.C},
\fpath{scripts/make_fit.C} (method 0) and \fpath{scripts/make_big_m2_tables.py}. In the overlay, \ovbigSEkeep\% of the SE-matched
tracks keep their own hit, \ovbigSEswap\% switch to a partner hit and \ovbigSElose\% lose the match (four-file values in
Sec.~\ref{sec:mudst}).

\subsubsection{Random matches: method 2 vs method 1}
Method 2 gives \ovbigFrCi\% random matches in 0--5\% and \ovbigFrAll\% in 0--80\% ($0.2<p<3$\,GeV/$c$). The two methods agree within
\bigTwoOneFrDev\% (relative) in every centrality group (Table~\ref{tab:bigm2cent}), although method 2 uses production
calibration and production matching while method 1 emulates both. Between the p and d peaks SE/ME is \ovbigHiCi{} (0--5\%), with
\bigTwoHiMidMin--\bigTwoHiMidMax{} per momentum bin at 0.6--2\,GeV/$c$ and \bigTwoHiLow{} at 0.2--0.4\,GeV/$c$
(Table~\ref{tab:bigm2bins}). The largest bin-by-bin ME/SE at 0.4--1.5\,GeV/$c$ is \ovbigVmaxCi{} (0--5\%), against
\mubigVmaxProdCi{} for method 1: method 2 has no timing artefact below the pion peak.

\input{tables/big_m2_cent}

\begin{figure}[H]
\centering
\figpanel{0.9}{ovbig_m2_cent}{$m^2$ of QA tracks, 0--5\%, nine momentum bins, 102 MuDst files (log scale). Black points: same event;
red filled: mixed event, method 2 (overlay with partner tracks, production makers re-run), normalized per QA track.
Random matches \ovbigFrCi\%; SE/ME between p and d \ovbigHiCi.}
\caption{Large sample, method 2, SE vs ME, 0--5\% (\fpath{scripts/make_ov.C}).}
\end{figure}

\begin{figure}[H]
\centering
\figpanel{0.49}{ovbig_m2_mid}{30--40\%: random matches \ovbigFrMid\%, SE/ME between p and d \ovbigHiMid.}
\figpanel{0.49}{ovbig_m2_per}{60--80\%: random matches \ovbigFrPer\%, SE/ME between p and d \ovbigHiPer.}
\caption{Large sample, method 2, SE vs ME in mid-central and peripheral events; encodings as in Fig.~\ref{fig:ovbig_m2_cent}.}
\end{figure}

\input{tables/big_m2_bins}

\begin{figure}[H]
\centering
\figpanel{0.9}{qm2_ovbig_cent}{$q\times m^2$ of QA tracks, 0--5\%, 102 MuDst files, 0.02\,GeV$^2$/$c^4$ bins (log scale). Black line:
same event; red filled: mixed event (method 2); blue points: SE$-$ME.}
\caption{Large sample, method 2, charge-signed $m^2$, 0--5\% (\fpath{scripts/make_qm2.C}).}
\end{figure}

\subsubsection{Signal extraction with method 2}
The fits of Sec.~\ref{sec:bigfit} were repeated with the method-2 ME (Fig.~\ref{fig:fit_ovbig_cent}, Table~\ref{tab:bigm2fit}). The
split Student-t (T2) is again best: in \bigTwoFitTtwoBestCi{} of 9 bins in 0--5\% and \bigTwoFitTtwoBestAll{} of 9 in 0--80\%. The
exception in both is 0.2--0.4\,GeV/$c$, where the double Gaussian is marginally better and no model is acceptable. Mean
$\chi^2$/ndf in 0--5\% is G2 \bigTwoFitChiGtwoCi, T \bigTwoFitChiTCi, T2 \bigTwoFitChiTtwoCi, lower than with method 1
(T2 \bigFitChiTtwoCi). The fitted yields add up to SE$-$ME within \bigTwoSumYDev\%. The yields of the two methods agree within
\bigTwoYieldAgree\% for $\pi$, K and p at 0.4--2\,GeV/$c$; at 0.2--0.4\,GeV/$c$ the method-2 proton yield is \bigTwoYieldLowP\% higher
(0--5\%), because method 2 subtracts less background there.

\choice{For the signal yields the choice between methods 1 and 2 matters only below 0.4\,GeV/$c$ and for K above 2\,GeV/$c$. Method 2
is the baseline (production calibration and matching); method 1 is a cross-check.}

Per charge (method 2, 0--5\%), the Student-t family is better in \bigTwoFitNegTfamCi{} (negative) and \bigTwoFitPosTfamCi{}
(positive) of 9 bins. The raw ratios at 0.4--2.5\,GeV/$c$ are $\pi^-/\pi^+ = $ \bigTwoRatPiMin--\bigTwoRatPiMax,
K$^-$/K$^+ = $ \bigTwoRatKMin--\bigTwoRatKMax{} and $\bar{\rm p}$/p $=$ \bigTwoRatPMin--\bigTwoRatPMax, the same as with method 1.

\begin{figure}[H]
\centering
\figpanel{0.9}{fit_ovbig_cent}{SE$-$ME $m^2$ (points), 0--5\%, all tracks, method 2 ME (log scale), with the signal part of the best
Gaussian-family model (blue dashed) and the best Student-t-family model (red).}
\caption{Large sample, method 2: SE$-$ME and line-shape fits, 0--5\% (\fpath{scripts/make_fit.C}).}
\end{figure}

\input{tables/big_m2_fit}

\subsubsection{The SE excess at low momentum}
\label{sec:lowp}
Both methods leave SE/ME $\approx2$ between the p and d peaks at 0.2--0.4\,GeV/$c$ in central events, and more in peripheral events.
\fpath{scripts/make_lowp.C} (log \fpath{logs/lowp_numbers.txt}) tests whether this excess behaves like random matches.

\paragraph{Centrality dependence.} Fig.~\ref{fig:lowp_cent} shows SE, ME and SE$-$ME in $1.6<m^2<2.8$\,GeV$^2$/$c^4$ per $10^4$
tracks (method 1: projected tracks; method 2: QA tracks) in the nine centrality classes. From 70--80\% to 0--5\% the ME per track rises
by a factor \lowpTwoAMErat{} (method 2, 0.2--0.4\,GeV/$c$), as random matches should. SE$-$ME stays flat:
\lowpTwoADmin--\lowpTwoADmax{} per $10^4$ tracks at 0.2--0.4\,GeV/$c$, \lowpTwoBDmin--\lowpTwoBDmax{} at 0.4--0.6 and
\lowpTwoCDmin--\lowpTwoCDmax{} at 0.6--1\,GeV/$c$ (method 2). Method 1 gives \lowpOneADmin--\lowpOneADmax,
\lowpOneBDmin--\lowpOneBDmax{} and \lowpOneCDmin--\lowpOneCDmax{} (per projected track). The excess is a fixed rate per track,
independent of the hit density. It is therefore not under-subtracted random background, which would scale like the ME. It stands
out at low momentum and in peripheral events only because the random background is small there.

\paragraph{Charge.} Table~\ref{tab:lowpcharge} gives the ratio of SE$-$ME of negative to positive tracks in
$1.6<m^2<2.8$\,GeV$^2$/$c^4$. Both methods agree. Two components are visible:
\begin{itemize}
\item at 0.3--0.6\,GeV/$c$ outside 60--80\% the ratio is \lowpChMidMin--\lowpChMidMax{} (method 2): most of the excess is
  charge-symmetric;
\item at 0.2--0.25\,GeV/$c$ the ratio drops to \lowpChLowMin--0.2: the excess is almost only positive. A natural candidate is slow
  protons that lose energy in the material before the TOF, arrive late and are reconstructed at larger $m^2$ (the high-$m^2$ tail of the
  proton peak that also requires the split Student-t). The ratio below 1 at 0.6--1\,GeV/$c$ is the same proton tail.
\end{itemize}
The charge-symmetric part excludes kaon decays in flight, $\Lambda$ decays and deuterons as its source.

\paragraph{Time structure.} Fig.~\ref{fig:lowp_dt} shows $t-t_\pi$ of SE and ME (method 2, 0--80\%). Beyond the particle peaks the SE$-$ME
forms a flat plateau up to the end of the window (27\,ns), at a level comparable to the ME. In 15--27\,ns the SE/ME is \lowpLateRatA{}
at 0.2--0.3\,GeV/$c$ and \lowpLateRatD{} at 0.6--1\,GeV/$c$, with a similar SE$-$ME per ns in all momentum bins
(\fpath{logs/lowp_numbers.txt}). There is no peak at a fixed delay. The $m^2$ window $1.6<m^2<2.8$ maps to delays of 15--35\,ns at
0.3\,GeV/$c$ but to a few ns at 1\,GeV/$c$, where the proton peak dominates; so a delay-flat component appears in the $m^2$ window
mainly at low momentum.

\input{tables/lowp_charge}

\caveat{Interpretation (not proven). The positive-only part at the lowest momenta is most likely the proton energy-loss tail. The
charge-symmetric part is a late hit correlated with the track itself, in or next to its cell, broad in time. Candidates are
secondaries from hadronic interactions of the particle in the TOF trays and the material behind them (slow nucleons, fragments,
back-splash), or delayed signals in neighbouring channels. A fixed-delay crosstalk is disfavoured by the flat time distribution. No ME can
model such a correlated component. If needed, it can be estimated as a per-track template, for instance from peripheral events where the
ME is negligible. Not done; see Sec.~\ref{sec:open}.}

\begin{figure}[H]
\centering
\figpanel{0.95}{lowp_cent}{SE (black), ME (red) and SE$-$ME (blue open) in $1.6<m^2<2.8$\,GeV$^2$/$c^4$ per $10^4$ tracks versus centrality,
for 0.2--0.4, 0.4--0.6 and 0.6--1\,GeV/$c$. Top: method 1 (per projected track); bottom: method 2 (per QA track).}
\caption{Large sample: centrality dependence of the SE excess between the p and d peaks (\fpath{scripts/make_lowp.C}).}
\end{figure}

\begin{figure}[H]
\centering
\figpanel{0.9}{lowp_dt}{$t-t_\pi$ of TOF-matched QA tracks, 0--80\%, four momentum bins (log scale). Black line: same event; red filled: mixed
event (method 2, normalized per momentum bin); blue points: SE$-$ME. Peaks: electrons just below 0, $\pi$ at 0, K and p at positive delays.}
\caption{Large sample, method 2: time structure of the SE excess (\fpath{scripts/make_lowp.C}).}
\end{figure}
